% Section 8: Proof of Theorem~\ref{ThmBdFinRank}

The goal of this section is to prove the theorems and the corollary in
the introduction. To this end let $g\ge 2$; we retain the notation of
$\mathsection$\ref{subsec:univfamily}. In particular,
$\pi\colon\mathfrak{A}_g\rightarrow \mathbb{A}_g$ is the universal
family of principally polarized abelian varieties of dimension $g$
with level-$\ell$-structure where $\ell\ge \bestlevel$ and
$\tau\colon \mathbb{M}_g\rightarrow \mathbb{A}_g$ is the Torelli morphism.

\begin{proposition}
  \label{prop:premazur}        
  The exist constants $c_1\ge 0,c_2\ge 1$ depending on the choices made
  above with the following property. Let $s\in \mathbb{M}_g(\IQbar)$
  with $h_{\overline{\mathbb{A}_g},\cM}(\tau(s)) \ge c_1$. Suppose $\Gamma$ is
  a finite rank subgroup of $\mathfrak{A}_{g,\tau(s)}(\IQbar)$ with rank
  $\rho\ge 0$. If $P_0 \in \mathfrak{C}_s(\IQbar)$, then
  $$
    \# (\mathfrak{C}_s(\IQbar)-P_0)\cap \Gamma \le
    c_2^{1+\rho}.
  $$
\end{proposition}


The proof combines Vojta's approach to the Mordell Conjecture with the
results obtained in $\mathsection$\ref{SectionDistanceCurve}. We
will use R\'emond's quantitative
version~\cite{Remond:Decompte,remond:vojtasup} of Vojta's method. A
similar approach was used in the authors's earlier work~\cite{DGH1p}
which also contains a review of Vojta's method in $\mathsection$2. Let
us recall the fundamental facts before proving Proposition~\ref{prop:premazur}.

Suppose we are given an abelian variety $A$ of dimension $g$ that is
defined over $\IQbar$ and is presented with a symmetric and very
ample line bundle $L$. We assume also that we have a closed immersion
of $A$ into some projective space $\IP_\IQbar^n$ determined by a basis
of the global sections of $L$. We assume that $A$ becomes a
projectively normal subvariety of $\IP_\IQbar^n$. This is the case if
$L$ is an at least fourth power of a symmetric and ample line bundle. 

Suppose $C$ is an irreducible  curve in $A$. Then let $\deg
C$ denote  the degree of $C$ considered as subvariety of
$A\subseteq \IP_\IQbar^n$, \textit{i.e.}, $\deg C = (C.L)$. Moreover, let $h(C)$ 
denote the height of $C$. 

On the ambient projective space  we have the 
 Weil height $h\colon
\IP_\IQbar^n(\IQbar)\rightarrow [0,\infty)$. Tate's Limit Argument,
compare \eqref{eq:fiberwiseNT},
applied to $h$ yields the N\'eron--Tate height $\hat h_L\colon A(\IQbar)\rightarrow
[0,\infty)$. It vanishes precisely on the points of finite order.
Moreover, it follows  from Tate's construction  that 
there exists a constant $c_{\mathrm{NT}}\ge 0$, which depends on $A$, such that
\begin{equation}
  \label{def:cNT}
  |\hat h_{L}(P) - h(P)|\le c_{\mathrm{NT}}
\end{equation}
for all $P\in A(\IQbar)$.

Finally, we need a measure for the heights of homogeneous polynomials that
define the addition and substraction on $A$, as required in R\'emond's
~\cite{remond:vojtasup}. Consider the $n+1$ global sections of $\cO(1)$
corresponding to the projective coordinates of $\mathbb{P}_\IQbar^n$.
They restrict to  global sections $\xi_0,\ldots,\xi_n$ of
$L$ on $A$. Let $f\colon A\times A\rightarrow A\times A$ denote the
morphism induced by $(P,Q)\mapsto (P+Q,P-Q)$, and let $p_1,p_2\colon
A\times A\rightarrow A$ be the first and section projection,
respectively. For all $i,j\in \{0,\ldots,n\}$ there are $P_{ij} \in
\IQbar[\mathbf{X},\mathbf{X}']$ with 
\begin{equation}
  \label{eq:Pijforh1}
  f^*(p_1^*\xi_i\otimes
  p_2^*\xi_j)=P_{ij}((p_1^*\xi_0,\ldots,p_1^*\xi_n),(p_2^*\xi_0,\ldots,p_2^*\xi_n))
\end{equation}
and where $P_{ij}$ is bihomogeneous of bidegree $(2,2)$ in $\mathbf{X}
= (X_0,\ldots,X_n)$ and $\mathbf{X}' = (X'_0,\ldots,X'_n)$; see 
 \cite[Proposition~5.2]{remond:vojtasup} with $a=b=1$
% $\mathsection$5.1 and $\mathsection$5.4~\cite{remond:vojtasup} 
for the existence of the $P_{ij}$. Here we require that
$\xi_0,\ldots,\xi_n$ constitute a basis of $H^0(A,L)$. Let $h_1$
denote the Weil height of the point in projective space whose
coordinates are all coefficients of all $P_{ij}$.

We point out a minor omission in \cite[$\mathsection$2]{DGH1p}: $h_1$
there must also involve both addition and subtraction on $A$, and not just
the addition.

The lemma below is \cite[Corollary~2.3]{DGH1p} which is itself a
standard application of R\'emond's explicit formulation of the Vojta
and Mumford inequalities. We thus obtain a bound that is exponential in the
rank of the subgroup $\Gamma$ for points of sufficiently large
N\'eron--Tate height. 

\begin{lemma}
  \label{lem:vojtamumford}
  Let $C$ be an irreducible curve in $A$.
There exists a constant $c=c(n,\deg C)\ge 1$ depending
only on $n$ and $\deg C$ with the following property. Suppose $\Gamma$
is a subgroup of $A(\IQbar)$ of finite rank $\rho\ge 0$. If $C$ is not
the translate of an algebraic subgroup of $A$, then
\begin{equation*}
\# \left \{P \in C(\IQbar)\cap \Gamma : \hat h_{L}(P) >
c\max\{1,h(C),c_{\mathrm{NT}},h_1\} \right\} \le c^\rho.
\end{equation*}
\end{lemma}

\begin{proof}[Proof of Proposition \ref{prop:premazur}]
  As in $\mathsection$\ref{subsec:univfamily} %we have finitely many
%  Zariski open affine subsets
%  $V_1,\ldots,V_t$ that cover $\mathbb{A}_g$. Moreover, 
  we have a
  closed immersion 
  $\mathfrak{A}_g \rightarrow\IP^n_\IQbar\times \mathbb{A}_g$ over $\mathbb{A}_g$.
  Let $s\in \mathbb{M}_g(\IQbar)$  with
  \begin{equation}
    \label{eq:hslbc1}
    h_{\overline{\mathbb{A}_g},\cM}(\tau(s)) \ge \max\{1,c_1\},
  \end{equation}
  where $c_1$
  comes from Proposition \ref{PropAlgPtFar} applied to $S = \mathbb{M}_g$.

 % Fix $i\in \{1,\ldots,t\}$ such that $\tau(s)\in V_i(\IQbar)$. 
  We now bound two quantities attached to the abelian variety
  $A=\mathfrak{A}_{g,\tau(s)}$ taken with its closed immersion into
  $\IP_\IQbar^n$. Observe that this closed immersion satisfies the
  condition imposed at the beginning of this section with $L=\cL|_A$
  where $\cL$ is as in $\mathsection$\ref{subsec:univfamily}.
  These quantities may depend on $s$. Below, $c>0$ denotes a constant
  that depends on the fixed data such as $g,n,$ and the ambient
  objects such as $\mathfrak{A}_g$ but not on  $s$.
  We will increase $c$ freely and without notice.

  \medskip
  \textbf{Bounding $c_{\mathrm{NT}}$.} For this we require the
  Silverman--Tate Theorem, Theorem \ref{thm:silvermantate}, applied to
  $\pi \colon \mathfrak{A}_g \rightarrow \mathbb{A}_g$. Recall that $h$ is the Weil 
  height on $\IP^n_{\IQbar}(\IQbar)$.
  % we may use it in lieu of
  % $h_{\overline{\cA},\overline{\cL}}$ in Theorem \ref{thm:silvermantate}.
  For all $P\in
  A(\IQbar)$ we have $|h(P)-\hat h(P)|\le c
  \max\{1,h_{\overline{\mathbb{A}_g},\cM}(\tau(s))\}$;
  note that we can bound $h_{\overline S}(\pi(P))$ from above linearly
  in terms of $h_{\overline{\mathbb{A}_g},\cM}(\tau(s))$ by the Height Machine.
  So we may take 
  \begin{equation}
    \label{eq:choosecNT}
    c_{\mathrm{NT}} = c\max\{1,h_{\overline{\mathbb{A}_g},\cM}(\tau(s))\}.    
  \end{equation}

  \medskip
  \textbf{Bounding $h_1$.}\label{page:boundh1} % Let $\xi_0,\ldots,\xi_n$ denote
  % global sections of $\cO(1)$ corresponding to 
  % projective coordinates of $\IP^n_{\IQbar}$.
  Recall that $f\colon A^2\rightarrow A^2$ sends $(P,Q)$ to
  $(P+Q,P-Q)$. We know that $P_{ij}$ as above exist. Here we will
  construct such a family with controlled height. To this end we
  consider points $P=[\zeta_0:\cdots:\zeta_n],Q=[\eta_0:\cdots:\eta_n]
  \in A(\IQbar)$. Then $f(P,Q) = ([\nu^+_0:\cdots
  :\nu^+_n],[\nu^-_0:\cdots:\nu^-_n])$. Recall that
  $A=\mathfrak{A}_{g,\tau(s)}$ is presented as a projectively normal
  subvariety of $\IP^n_\IQbar$ by the construction in
  $\mathsection$\ref{subsec:univfamily}. By \eqref{eq:Pijforh1}  there
  is for each $i,j\in \{0,\ldots,n\}$
  a bihomogeneous polynomial $P_{ij}$ of bidegree $(2,2)$ that is
  independent of $P$ and $Q$, with
  \begin{equation}
    \label{eq:inhomlinearsystemforh1}
    \nu^+_i \nu^-_j = \lambda P_{ij}((\zeta_0,\ldots,\zeta_n),(\eta_0,\ldots,\eta_n))
  \end{equation}
  for some non-zero $\lambda\in\IQbar$ that may depend on $(P,Q)$. 
  We  eliminate $\lambda$ and consider  
  \begin{equation}
    \label{eq:linearsystemforh1}
    \nu^+_i \nu^-_jP_{i'j'}((\zeta_0,\ldots,\zeta_n),(\eta_0,\ldots,\eta_n)) - \nu^+_{i'} \nu^-_{j'}P_{ij}((\zeta_0,\ldots,\zeta_n),(\eta_0,\ldots,\eta_n))=0
  \end{equation}
  as a system of homogeneous linear equations 
  parametrized by $(i,j),(i',j')\in \{0,\ldots,n\}^2$, the unknowns are
  the coefficients of the $P_{ij}$. As each $P_{ij}$ is bihomogeneous
  of bidegree $(2,2)$, the number of unknowns is
  $N=(n+1)^2 {n+2 \choose 2}^2$ which is independent of $s$. 

  Each pair of points $(P,Q)\in A(\IQbar)^2$ yields one system of
  linear equations.
  We know that there is a non-trivial solution $(P_{ij})_{ij}$ that solves for all
  $(P,Q)$ simultaneously and such that some $P_{ij}$ does not vanish
  identically on $A\times A$. Our goal is to find such a common solution of
  controlled height.

  First, observe that a common solution  for when
  $(P,Q)$ runs over all torsion points of $A^2(\IQbar)$ is a common
  solution for all pairs $(P,Q)$. Indeed, this follows as torsion
  points of $A^2(\IQbar)$ lie Zariski dense in $A^2$. 
  Second, observe that the full system has finite rank $M< N$ so it
  suffices to consider only finite many torsion points $(P,Q)$.

  Our task is thus to find a  common solution to
  (\ref{eq:linearsystemforh1}) for all $(i,j),(i',j')$, where some
  $P_{ij}$ does not vanish identically on $A\times A$, and where 
  $[\zeta_0:\cdots:\zeta_n]$,$[\eta_0:\cdots:\eta_n]$, and
  $[\nu^{\pm}_0:\cdots :\nu^{\pm}_n]$ are certain torsion points on
  $A(\IQbar)$. We may assume that some $\zeta_i$ is $1$ and similarly
  for  $\eta_i$ and
  $\nu^{\pm}_i$. So all coordinates are in $\IQbar$. Moreover, the
  height  of each torsion point is at most $c_{\mathrm{NT}}$ by (\ref{def:cNT}).
  The resulting system of linear equations is represented by an
  $M\times N$ matrix with algebraic coefficients. By elementary
  properties of the height, each coefficient in the system has affine
  Weil height $c \cdot c_{\mathrm{NT}}$, for $c$ large enough. It is tempting, but
  unnecessary, to invoke Siegel's Lemma to find a non-trivial solution.
  As $M,N$ are bounded in terms of $n$, 
  Cramer's Rule establishes the existence of a basis of 
  non-zero solution such that the Weil height of the coefficient
  vector is at most $c \cdot c_{\mathrm{NT}}$.
  Among this basis there
  is one solution where one $P_{ij}$ does not vanish identically on
  $A\times A$. By \eqref{eq:choosecNT} we find
  \begin{equation}
    \label{eq:h1bound}
    h_1 \le     c\max\{1,h_{\overline{\mathbb{A}_g},\cM}(\tau(s))\}
  \end{equation}
  for the projective height of the tuple $(P_{ij})_{ij}$. 
  Thus \eqref{eq:linearsystemforh1} holds and so we get
  \eqref{eq:inhomlinearsystemforh1} on all of $A^2(\IQbar)$, at least with  $\lambda$ a rational
  function on $A\times A$ that is not identically zero. But $\lambda$ cannot
  vanish anywhere on $A\times A$, as otherwise the left-hand side of
  \eqref{eq:inhomlinearsystemforh1} would vanish at some point of
  $A\times A$ for all $(i,j)$. Hence $\lambda$ is a non-zero constant. Replacing
  $P_{ij}$ by $\lambda P_{ij}$ does not change the projective height;
  we get \eqref{eq:Pijforh1} with the desired bound for $h_1$. 

  \medskip
  \textbf{Bounding height and degree of a curve.} 
  By Lemma \ref{lem:uniformdegreebound} we have
  \begin{equation}
    \label{eq:degreehgtbound}
    \deg(\mathfrak{C}_s-P_s)\le c\quad\text{and}\quad 
    h(\mathfrak{C}_s-P_s)\le
    c \max\{1,h_{\overline{\mathbb{A}_g},\cM}(\tau(s))\}
  \end{equation}
  for some
  $P_s\in \mathfrak{C}_s(\IQbar)$.

  \bigskip

  We now follow the argumentation in~\cite{DGH1p}. Let $\Gamma$ be a
  subgroup of $\mathfrak{A}_{g,\tau(s)}(\IQbar)$ for finite rank $\rho$.
  We first prove the proposition in the case $P_0=P_s$. We apply Lemma
  \ref{lem:vojtamumford}  to the curve $C =
  \mathfrak{C}_s-P_s\subset \mathfrak{A}_{g,\tau(s)}=A$ and use the bounds
  (\ref{eq:choosecNT}), \eqref{eq:h1bound}, and 
  (\ref{eq:degreehgtbound}).
  Note that $C$ is a smooth curve of genus
  $g\ge 2$. So it cannot be the translate of an algebraic subgroup of
  $A$.
  It follows that the number of points $P\in
  \mathfrak{C}_s(\IQbar)$ with $P-P_s\in \Gamma$ and $\hat h(P-P_s) >
  R^2$ where 
  \begin{equation}
    \label{def:bigR}
    R= (c\max\{1,h_{\overline{\mathbb{A}_g},\cM}(\tau(s))\})^{1/2}    
  \end{equation}
  is at most  $c^{\rho}\le c^{1+\rho}$.

  The burden of this paper is  to find a bound of the same quality 
  for the number of pairwise distinct points $P_1,P_2,P_3,\ldots$  in
  $\mathfrak{C}_s(\IQbar)$ with $\hat h(P_i-P_s) \le R^2$. This is where
  Proposition~\ref{PropAlgPtFar} enters. Recall our assumption
  (\ref{eq:hslbc1}) on $s$.
  Let $c_2'$ be the constant $c_2$ from Proposition~\ref{PropAlgPtFar}; it is independent of $s$. As
  $\# \Xi_s\le c_2'$ we may assume $P_i\not\in \Xi_s$ for all $i$. 
  So we may assume that each $P_i$ is in the second alternative of
  Proposition~\ref{PropAlgPtFar}.

  As in \cite{DGH1p} we use the Euclidean norm $|\cdot|$ defined by $\hat h^{1/2}$ on
  the $\rho$-dimensional  $\IR$-vector space $\Gamma\otimes\IR$. 
  Let $r\in (0,R]$. By an elementary ball packing argument, any subset of
  $\Gamma\otimes\IR$ contained in a closed ball of radius $R$ is
  covered by at most $(1+2R/r)^{\rho}$ closed balls of radius $r$ centered at
  elements of the given subset; see \cite[Lemme~6.1]{Remond:Decompte}.
  We apply this geometric argument to $R$ as in (\ref{def:bigR}) and
  to  $r$, the positive square-root of 
  $h_{\overline{\mathbb{A}_g},\cM}(\tau(s))/c_3 =
  \max\{1,h_{\overline{\mathbb{A}_g},\cM}(\tau(s))\}/c_3$.
  The contribution of the height
  $h_{\overline{\mathbb{A}_g},\cM}(\tau(s))$ cancels in the quotient
  $R/r$. We find $R/r\le c$. So the
  number of balls in the covering is at most $c^{1+\rho}$. 

  By Proposition~\ref{PropAlgPtFar}(ii) the number of the $P_i$'s that
  map to a single closed ball of radius $r$ is
  at most $c_4$. Thus after increasing $c$ we find that $\#\{P_i \in \mathfrak{C}_s(\IQbar)\cap \Gamma : \hat{h}(P_i-P_s) \le R^2\} \le c_4 c^{1+\rho}$, %the
%  number of $P_i$ as above is at most $c_4 c^{1+\rho}$, 
as desired. 
  This completes the proof of the proposition in the case $P_0=P_s$
  for sufficiently large $c_2$.

  The case of a general base point follows easily as our estimates
  depend only on the rank $\rho$ of $\Gamma$. Indeed,
  let $P_0\in \mathfrak{C}_s(\IQbar)$ be an arbitrary point and let
  $\Gamma'$ be the subgroup of $\mathfrak{A}_{g,\tau(s)}(\IQbar)$ generated by
  $\Gamma$ and $P_0-P_s$. Its rank is at most $\rho +1$.

  Now if $Q\in \mathfrak{C}_s(\IQbar)-P_0$ lies in $\Gamma$, then
  $Q+P_0-P_s\in \mathfrak{C}_s(\IQbar)-P_s$ lies in $\Gamma'$. The
  number of such $Q$ is at most $c_2^{2+\rho}$ by
  what we already proved. 
  The proposition follows as $c_2^{2+\rho}\le (c_2^2)^{1+\rho}$ and
  since we may replace $c_2$ by $c_2^2$. 
\end{proof}




\begin{proof}[Proof of Theorem~\ref{ThmBdRatIntro}] 
  It is possible to deduce Theorem~\ref{ThmBdRatIntro} from
  Theorem~\ref{ThmBdFinRank}, which we prove below. However in view of
  the importance of Theorem~\ref{ThmBdRatIntro}, we hereby give it a
  complete proof.

  This proof works for any level $\ell\ge\bestlevel$, but we may fix
  $\ell=\bestlevel$ for definiteness. Let
  $\mathbb{A}_g,\overline{\mathbb{A}_g},\cM,$
  and $h_{\overline{\mathbb{A}_g},\cM}$ be as in $\mathsection$\ref{subsec:univfamily}.

  
  Our curve $C$ corresponds to an $F$-rational point $s_F$ of $\mathbb{M}_{g,1}$, the coarse
  moduli space of smooth genus $g$ curves without level structure. 

  
  The fine moduli space $\mathbb{M}_g$ of smooth genus $g$ curves
  with level-$\ell$-structure is a finite cover of $\mathbb{M}_{g,1}$.
  For this proof it is convenient to recall that 
  $\mathbb{M}_g$ is defined over
  the cyclotomic field generated by a third root of unity;
  recall the convention that we fixed a third root of unity and that
  $\mathbb{M}_g$ is geometrically irreducible.
  % with $\zeta$ a root of unity of order $\ell$. 
  %After replacing  $F$ by a finite extension whose degree is bounded solely in terms of $g$ we may assume that $\mathbb{M}_g$ is defined over $F$. 
%  Moreover, there exists a finite extension $F'/F$
%  with $F'\subset \IQbar$ and $s\in \mathbb{M}_g(F')$
  Say $s\in\mathbb{M}_g(\IQbar)$ maps to $s_F$. Then
  $F'=F(s)$ is a number field and  
  % Thus there exist a number field $F'\subset\IQbar$ containing $F$
  % % finite extension $F'/F$, with $F'\subset\IQbar$, and  
  % and $s \in \mathbb{M}_g(F')$ 
  % that maps to $s_F$. Moreover, we may assume that
  $[F':F]$ is bounded above
  only in terms of $g$ and $\ell$. We may identify $C_{F'} = C\otimes_F F'$ with $\mathfrak{C}_s$, the 
  fiber of
  $\mathfrak{C}_g\rightarrow\mathbb{M}_g$ above $s$.
  %Thus $C(K)\subset C(L) = C_L(L) = \mathfrak{C}_s(L)$.

  Constructing the
  Jacobian commutes with finite field extension. We thus view 
  $\Gamma = \mathrm{Jac}(C)(F)$ as a subgroup of $\mathrm{Jac}(C)(\IQbar)
  = \mathrm{Jac}(C_{F'})(\IQbar)$. 


  To prove the theorem we may assume $C(F)\not=\emptyset$. So fix
  $P_0\in C(F)$. We consider the Abel--Jacobi embedding
  $C-P_0\subset \mathrm{Jac}(C)$ defined over $F$. Then $\# C(F)\le \#
  (C_{F'}(\IQbar)-P_0)\cap\Gamma = \#
  (\mathfrak{C}_s(\IQbar)-P_0)\cap \Gamma$. If
  $h_{\overline{\mathbb{A}_g},\cM}(\tau(s))\ge c_1$, the theorem
  follows from Proposition~\ref{prop:premazur}. 
  Note that in this case, the constant $c$ in \eqref{eq:CKexpinrank} is independent of $d$.

  So we may assume that the height 
  of $\tau(s)$ is less than $c_1$. 
  As $[F':\IQ]\le [F':F][F:\IQ]\le [F':F]d$,
  Northcott's Theorem implies
 that $\tau(s)$ comes   from a finite set in $\mathbb{A}_{g}(\IQbar)$
 that depends only on $g,d,$ and $\ell$. The same holds for $s$ and
 thus $F'=F(s)$
 since the Torelli morphism $\tau$ is finite-to-$1$ and thus   has
 fibers of bounded cardinality. % We may also assume that $F'$ comes
 % from a finite set. 
 This means that the remaining $C$ are twists   in
 finitely many $F'$-isomorphism classes. But then it suffices to
 apply R\'emond's estimate \cite[page 643]{DPvarabII} to a single
 $C_{F'}$ and    use $\#C(F) \le \# (C_{F'}(\IQbar)-P_0)\cap\Gamma$ to
 conclude the theorem. % as we have fixed $\ell = 3$.
% \Pcomment{The commment that $c$ is polynomial in
 %  $d$ is now in remark 8.3. Why not keep it here in the proof and 
% add justifications as desired by
 %  the referee?} \Ginlinecomment{My impression is that we have to repeat the last paragraph of the proof of Thm~1.2, in addition to %the current argument of Remark~8.3, for this. Plus, it would be nice to explain why Thm~1.2 implies Thm~1.1.}

% \Pinlinecomment{I put the reference to Silverman here because it doesn't really fit that   well after the proof any.}
Silverman's older result~\cite[Theorem 1]{Silverman:twists} also
 handles uniformity among twists.
\end{proof}

%\Pinlinecomment{New arguments start here.}
Let us explain how to obtain some extra uniformity in the second case
of the proof. More precisely, we show that the constant $c(g,d)$ in
\eqref{eq:CKexpinrank} grows polynomially in $d$. We retain the above
proof's notation.% \Gcomment{Slightly rephrased to point out here that
%$c(g,d)$ depends polynomially on $d$.}\Pcomment{OK, but I think we can omit  ``when
%$h_{\overline{\mathbb{A}_g},\cM}(\tau(s)) < c_1$''}

Denote by $\rho \colon \mathbb{A}_g = \mathbb{A}_{g,\ell} \rightarrow
\mathbb{A}_{g,1}$ the natural morphism to the coarse moduli space which forgets the level structure.
We recall that
$\overline{\mathbb{A}_g}$ %and $\overline{\mathbb{A}_{g,1}}$
is presented as a closed subvariety
of
projective space induced by a basis of global sections of a positive
powers of ample line bundle $\cM$. Let
$\iota\colon\mathbb{A}_{g,1}\rightarrow\IP_{\mathbb{Q}}^m$ be an
immersion, as before Theorem~\ref{ThmBdFinRank}.
By 
\cite[Lemma~4]{Silverman:heightest11}, there exist $c' > 0 $ and $c''
\ge 0$ depending on the immersions such that
$h(\iota(\rho(t))) \le c'
h_{\overline{\mathbb{A}_g},\cM}(t) + c''$ for all $t\in
\mathbb{A}_g(\IQbar)$.
So $h(\iota(\rho(\tau(s)))) \le c'c_1+c''$
is bounded uniformly in the second case.

By fundamental work of Faltings~\cite[$\mathsection 3$ including
 the proof of Lemma~3]{Faltings:ES}, see also
\cite[the remarks below Proposition~V.4.4 and Proposition~V.4.5]{FaltingsChai}, the \textit{stable Faltings height} of
$\mathfrak{A}_{g,\tau(s)}$ is bounded from above in terms of
$c'c_1+c''$ and $g$ only.
The height $h_{\mathrm{DP}}(\mathfrak{A}_{g,\tau(s)})$ used by David and
Philippon is bounded similarly by work of Bost and David,
see~\cite[Corollaire~6.9]{DPvarabII} and \cite{Pazuki:12}. 

% \Pcomment{Reference for this would be nice.} \Ginlinecomment{I'm afraid that we have to repeat the last paragraph of the proof of Thm~1.2 for this.}
% that the
% stable Faltings height of $\mathfrak{A}_{g,\tau(s)}$ is bounded from
% above in terms of $c_1$ and $g$. The height
% $h(\mathfrak{A}_{g,\tau(s)})$ used in David and Philippon's paper is
% thus also bounded in terms of $c_1$, \textit{cf.}
% Corollaire 6.9~\cite{DPvarabII}.
In R\'emond's bound  \cite[page 643]{DPvarabII}
for $\#(C_{F'}(\IQbar)-P_0)\cap \Gamma$,
the base in the exponential depends polynomially on 
$D\max\{1,h_{\mathrm{DP}}(\mathfrak{A}_{g,\tau(s)})\}$,
where $D$ is the degree over
$\IQ$ of a suitable field of definition of $\mathfrak{A}_{g,\tau(s)}$.
As this abelian variety can be defined over $F'$ we may assume $D\le
[F':F]d$ is bounded linearly in $d$.
Recall that $\deg (C_{F'}-P_0)$ is bounded from above uniformly. 
% Keeping the above proof's notation.
So R\'emond's bound implies that  $c(g,d)$
in (\ref{eq:CKexpinrank}) can be chosen to grow at most
polynomially in $d$.

The definition of $h_{\mathrm{DP}}(\mathfrak{A}_{g,\tau(s)})$
involves  theta functions and a different kind of level structure.
Using standard results on heights and by
going down and up in the level structure it is likely that one can
bound $h_{\mathrm{DP}}(\mathfrak{A}_{g,\tau(s)})$ from above directly in terms of
$h_{\overline{\mathbb{A}_g},\cM}(\tau(s))$. For this one would need to
work with a different level $\ell$ in the proof of
Theorem~\ref{ThmBdRatIntro}.

\begin{proof}[Proof of Theorem~\ref{ThmBdFinRank}] 
  We keep the same notation as in the proof of
  Theorem~\ref{ThmBdRatIntro}. So $\ell = \bestlevel$ and
  $\mathbb{A}_g,\overline{\mathbb{A}_g},\cM,$ and
  $h_{\overline{\mathbb{A}_g},\cM}$
  are as in $\mathsection$\ref{subsec:univfamily}.

  Let $C$ be a smooth curve of genus $g\ge 2$ defined over $\IQbar$,
  and let $\Gamma$ be a finite rank subgroup of
  $\mathrm{Jac}(C)(\IQbar)$. Let $P_0 \in C(\IQbar)$.
  
  The curve $C$
  corresponds to a $\IQbar$-point $s_{\mathrm{c}}$ of
  $\mathbb{M}_{g,1}$.% \st{the coarse
%    moduli space of smooth genus $g$ curves without level structure.}
 % \Pinlinecomment{Propose to omit as it repeats what we said just a
  %  few lines further up.}
  
  
  The fine moduli space $\mathbb{M}_g$ of smooth genus $g$ curves
  with level-$\ell$-structure is a finite covering of $\mathbb{M}_{g,1}$. So there exists an $s \in \mathbb{M}_g(\IQbar)$ 
  that maps to $s_{\mathrm{c}}$. Thus $C$ is isomorphic, over $\IQbar$, to the fiber $\mathfrak{C}_s$ of the universal family 
  $\mathfrak{C}_g\rightarrow\mathbb{M}_g$. We thus view $\Gamma$ as a
  finite rank subgroup of $\mathrm{Jac}(\mathfrak{C}_s)(\IQbar)$, and
  $P_0 \in \mathfrak{C}_s(\IQbar)$. 
  
  Consider the Abel--Jacobi embedding $C-P_0 \subseteq
  \mathrm{Jac}(C)$. Then $\# (C(\IQbar)-P_0) \cap \Gamma = \#
  (\mathfrak{C}_s(\IQbar) - P_0)\cap \Gamma$. If $h_{\overline{\mathbb{A}_g},\cM}(\tau(s)) \ge c_1$, then $\# (C(\IQbar)-P_0) \cap \Gamma \le c_2^{1+\rho}$ 
  by Proposition~\ref{prop:premazur}. % by taking $\delta = c_1$ and $c = c_2$.

 % \Pcomment{I simplified this a bit because the same argument appears    after the proof of Theorem \ref{ThmBdRatIntro}.}
  Thus it suffices to find a constant $c'_1\ge 0$ that is independent
  of $C$ and such % = c'_1(g,\tau) \ge 0$ such
  that $h(\iota([\mathrm{Jac}(C)])) \ge c_1'$ implies
  $h_{\overline{\mathbb{A}_g},\cM}(\tau(s)) \ge c_1$. 
  
  As after the proof of Theorem~\ref{ThmBdRatIntro}
  denote by $\rho \colon \mathbb{A}_g = \mathbb{A}_{g,\ell}
  \rightarrow \mathbb{A}_{g,1}$ the natural morphism.
  % We assume to have presented immersions of
  % $\overline{\mathbb{A}_g}$ and $\mathbb{A}_{g,1}$
  % into projective spaces, the former induced by a basis of global
  % sections of  a positive power of
  % $\cM$.  
  As in the proof of Theorem~\ref{ThmBdFinRank} we use
  $h(\iota(\rho(t)))\le c'h_{\overline{\IA_{g}},\cM}(t)  + c''$ for all $t\in
  \mathbb{A}_g(\IQbar)$.
  The theorem follows since $\rho(\tau(s)) = [\mathrm{Jac}(C)]$.
  % and as $h(\tau(s))$ is from above (and below)
  % bounded linearly in terms
  % of $h_{\overline{\mathbb{A}_g},\cM}(\tau(s))$.
  % We can write $c' = c'(g,\tau)$ and $c''=c''(g,\tau)$. 
  %
%it suffices to take $c'_1 = (c_1+c'')/c'$.
  %we have $h(\rho(\tau(s))) \le \delta'$ for some $\delta' = \delta'(g)$. But $\rho(\tau(s)) = [\jac{C}]$. Hence $h_{\mathrm{Fal}}(\jac{C}) \le \delta$ for some $\delta = \delta(g)$ by Faltings; see \cite[]{}. Thus the theorem for this case follows immediately from R\'emond's estimate \cite[page 643]{DPvarabII}.    Note that in this case, the constant $c$ in \eqref{eq:CKexpinrank} is polynomial in $d$.
\end{proof}
\begin{remark}
It is possible to prove Theorem~\ref{ThmBdRatIntro} (without the
dependency claims on $c(g,d)$) using
Theorem~\ref{ThmBdFinRank}. Let $C$ be a smooth curve of genus $g \ge
2$ defined over a number field $F \subseteq \IQbar$. Then by taking $\Gamma =
\mathrm{Jac}(C)(F)$ in Theorem~\ref{ThmBdFinRank}, we can conclude
Theorem~\ref{ThmBdRatIntro} if $h(\iota([\mathrm{Jac}(C)])) \ge c_1$. 
%When
%$h(\iota([\mathrm{Jac}(C)])) \le c_1$, then
%$h_{\mathrm{Fal}}(\mathrm{Jac}(C)) \le \delta$ for some $\delta =
%\delta(g,c_1)$ by Faltings \cite[]{}.\Pcomment{Added dependency in
%  $c_1$.} Thus R\'emond's estimate \cite[page 643]{DPvarabII} yields
%the desired bound $c^{1+\mathrm{rank}\mathrm{Jac}(C)}$, and the
%constant $c$ is polynomial in $d$.
%\Pinlinecomment{Polynomial dependency in $d$ is now proved just after
%  the proof of thm 1.1.}\Pinlinecomment{I propose to replace the two
%  final sentences in this
%  remark  by something along the lines of "
  The case $h(\iota([\mathrm{Jac}(C)])) < c_1$ can be 
  handled as in the
  proof of Theorem~\ref{ThmBdRatIntro}, and one can furthermore obtain 
  extra uniformity
  for $c_2$ in Theorem~\ref{ThmBdFinRank} by applying R\'emond's bound
  \cite[page~643]{DPvarabII} as after the proof of Theorem~\ref{ThmBdRatIntro}.%". There's no need to explain the same
%  thing twice.}
\end{remark}

% \begin{proof}[Proof of Corollary~\ref{CorBdAlgTorIntro}]
%%% Requires no proof, as proof is explained in introduction
% \end{proof}

\begin{proof}[Proof of Theorem~\ref{thm:BdTorIntroNF}]
Let $C$ be a smooth curve of genus $g\ge 2$  defined over a number field $F\subset\IQbar$.
  
Apply Theorem~\ref{ThmBdFinRank} to $C_{\IQbar}$, $P_0 \in C(\IQbar)$
and $\Gamma = \jac(C)({\IQbar})_{\mathrm{tors}}$, whose rank is $0$. Then we obtain $c_1 \ge 0$ and $c_2 \ge 1$ such that
  \[
  \#(C(\IQbar) - P_0) \cap \jac(C)({\IQbar})_{\mathrm{tors}} \le c_2
  \]
  if $h(\iota([\mathrm{Jac}(C_\IQbar)])) \ge c_1$.
  
By the Northcott property and Torelli's Theorem, there are up-to
$\IQbar$-isomophism only finitely many $C_{\IQbar}$'s  defined over a number field $F$ with $[F:\IQ] \le d$ such that  $h(\iota([\mathrm{Jac}(C_\IQbar)])) < c_1$. By applying Raynaud's result on the Manin--Mumford Conjecture to each one of these finitely many curves separately, we obtain Theorem~\ref{thm:BdTorIntroNF}.
\end{proof}

%%% Local Variables:
%%% TeX-master: "main"
%%% End:
